\chapter{Caso documentado del TFG}
\label{app:tfg}

Este apéndice resume el montaje, el entorno de software y los resultados utilizados
durante el TFG. Corresponden a una configuración concreta y no deben interpretarse como
requisitos generales de la AKD1000.

\section{Montaje utilizado}

El montaje empleado durante el TFG estaba formado por los componentes recogidos en el
\cref{tab:componentes-tfg} \cite{poveda_tfg_2026}.

\begin{longtable}{@{}L{0.27\textwidth}L{0.43\textwidth}L{0.22\textwidth}@{}}
\caption{Hardware utilizado durante el TFG.}
\label{tab:componentes-tfg}\\
\toprule
\textbf{Elemento} & \textbf{Identificación} & \textbf{Observaciones} \\
\midrule
\endfirsthead

\toprule
\textbf{Elemento} & \textbf{Identificación} & \textbf{Observaciones} \\
\midrule
\endhead

Host &
Raspberry Pi 5 Model B, 8~GB de RAM, \code{Rev 1.1}. &
Adquirida dentro de un kit db-tronic. \\

Almacenamiento &
SSD NVMe M.2 de 64~GB. &
Incluido en el kit; no se conserva el modelo exacto. \\

Refrigeración y carcasa &
Sistema de refrigeración activa y carcasa incluidos en el kit. &
No se conserva la referencia comercial. \\

Fuente del host &
Fuente de alimentación de 27~W. &
Incluida en el kit. \\

Adaptador &
Waveshare \emph{PCIe TO PCIe x1 Board (C)} para Raspberry Pi 5. &
Conectado a la Raspberry Pi mediante FFC. \\

Cable &
FFC con serigrafía \code{PCIe-Cable-50mm}. &
Longitud de 50~mm. \\

Tarjeta &
BrainChip Akida PCIe Board AKD1000. &
El dispositivo utilizado aparecía como \code{BC.00.000.002}. \\

Alimentación de la tarjeta &
12~V en continua mediante el carrier. &
Alimentación independiente de la Raspberry Pi. \\

\bottomrule
\end{longtable}

La cadena física utilizada fue:

\begin{center}
Raspberry Pi 5
$\rightarrow$
cable FFC
$\rightarrow$
carrier Waveshare
$\rightarrow$
BrainChip AKD1000
\end{center}

La Raspberry Pi se alimentó mediante USB-C y la tarjeta AKD1000 mediante la entrada de
12~V del carrier.

\section{Entorno de software}

Las versiones utilizadas durante la campaña final se muestran en el
\cref{tab:tfg-env}.

\begin{table}[H]
\centering
\caption{Entorno de software del TFG.}
\label{tab:tfg-env}
\begin{tabularx}{\textwidth}{@{}L{0.32\textwidth}Y@{}}
\toprule
\textbf{Elemento} & \textbf{Versión} \\
\midrule
Sistema operativo & Debian GNU/Linux 12 (bookworm) \\
Python & 3.11.2 \\
TensorFlow & 2.19.1 \\
API Keras & \code{tf\_keras} \\
Akida & 2.19.1 \\
CNN2SNN & 2.19.1 \\
QuantizeML & 1.2.3 \\
Dispositivo & \code{BC.00.000.002} \\
Configuración del host & governor \code{performance}; \code{throttled=0x0} \\
\bottomrule
\end{tabularx}
\end{table}

Estas versiones corresponden al entorno utilizado en el experimento. Para reproducir
completamente los resultados también son necesarios los modelos, los datos y el mismo
procedimiento de medida.

\section{Modelos AoA}

El TFG utilizó dos redes para estimar el ángulo de llegada (\emph{Angle of Arrival}, AoA).
Cada muestra contiene ocho valores enteros de intensidad de envolvente entre 0 y 15.

La salida del modelo contiene 180 valores, uno por cada clase. La clase $c$ corresponde
al ángulo

\begin{equation}
\theta = 2c.
\end{equation}

Los dos modelos finales fueron:

\begin{table}[H]
\centering
\caption{Modelos AoA utilizados en el TFG.}
\label{tab:tfg-models}
\begin{tabularx}{\textwidth}{@{}L{0.23\textwidth}L{0.28\textwidth}L{0.23\textwidth}Y@{}}
\toprule
\textbf{Modelo} & \textbf{Arquitectura} & \textbf{Pesos sin sesgos} & \textbf{Uso} \\
\midrule
Compacto &
$8\rightarrow10\rightarrow180$ &
1.880 &
Modelo de baja complejidad \\

Principal &
$8\rightarrow957\rightarrow180$ &
179.916 &
Modelo con mayor precisión \\
\bottomrule
\end{tabularx}
\end{table}

El error angular se calculó teniendo en cuenta la naturaleza circular del ángulo:

\begin{equation}
e_{\mathrm{circ}}
=
((\hat{\theta}-\theta+180)\bmod 360)-180.
\end{equation}

Por ejemplo, una predicción de $0^\circ$ para una referencia de $358^\circ$ produce
un error de $2^\circ$ y no de $-358^\circ$.

Los modelos FBZ y el dataset utilizados en la campaña original no forman parte de esta
entrega. El modelo de prueba incluido con el repositorio se utiliza únicamente para comprobar
el funcionamiento del entorno y de la tarjeta.

\section{Resultados}

La evaluación se realizó sobre 576 muestras. El modelo principal obtuvo un MAE angular
de $4{,}09^\circ$, frente a $9{,}13^\circ$ del modelo compacto
\cite{poveda_tfg_2026}.

En las configuraciones estudiadas, los puntos de menor energía estimada fueron:

\begin{itemize}
  \item modelo compacto: \code{Economy + HwPr}, con 0,565~ms,
        770,0~mW y 435,2~$\mu$J;

  \item modelo principal: \code{Performance + AllNps}, con 0,756~ms,
        731,8~mW y 553,3~$\mu$J.
\end{itemize}

\begin{figure}[H]
  \centering
  \includegraphics[width=0.88\textwidth]{nine-mode-results.png}
  \caption[Resultados de los modelos AoA]
  {Latencia y energía estimada para los dos modelos en las nueve combinaciones de modo
  de reloj y mapeo evaluadas durante el TFG.}
  \label{fig:tfg-results}
\end{figure}

La latencia se midió desde el host alrededor de la llamada de ejecución. La potencia
corresponde a la medida proporcionada por la API de Akida y no al consumo completo de
la Raspberry Pi, el carrier y la tarjeta.

Estos resultados describen únicamente la plataforma utilizada durante el TFG.